\section*{Appendix A: Evaluation Schemas}
\addcontentsline{toc}{section}{Appendix A: Evaluation Schemas}
\label{sec:appendix-schemas}

A recurring challenge with existing document extraction benchmarks is incomplete or ambiguous schema documentation. Without clear specifications, it can be difficult to understand expected field formats, handling of optional values, or normalization rules. Researchers attempting to reproduce results must often reverse-engineer these details from examples or evaluation scripts, leading to inconsistent implementations and incomparable metrics.

To improve reproducibility, we provide explicit schemas for claim incidents and the main heterogeneous row families that benefit most from a public field contract. Other operations families are defined by their ground-truth field contracts and scored by the generic record-list evaluator. The evaluation scripts normalize model outputs before scoring so format errors are caught early rather than silently degrading metrics.

\subsection*{A.1 Financial Breakdown Schema}

Each incident contains four financial breakdown objects (\texttt{bi}, \texttt{pd}, \texttt{lae}, \texttt{ded}) with the following fields:

\begingroup
\small
\setlength{\tabcolsep}{3pt}
\renewcommand{\arraystretch}{1.15}
\begin{longtable}{@{}>{\raggedright\arraybackslash}p{0.19\textwidth} >{\raggedright\arraybackslash}p{0.13\textwidth} >{\raggedright\arraybackslash}p{0.13\textwidth} >{\raggedright\arraybackslash}p{0.43\textwidth}@{}}
\toprule
Field & Type & Default & Description \\
\midrule
\endfirsthead
\toprule
Field & Type & Default & Description \\
\midrule
\endhead
\midrule
\endfoot
\bottomrule
\endlastfoot
\texttt{reserve} & float & 0.0 & Amount reserved for potential payout \\
\texttt{paid} & float & 0.0 & Amount already paid \\
\texttt{recovered} & float & 0.0 & Amount recovered (e.g., deductible) \\
\texttt{total\_\allowbreak incurred} & float & 0.0 & Reserve + Paid - Recovered \\
\end{longtable}
\endgroup

\subsection*{A.2 Loss Run Incident Schema}

The primary entity schema representing a single insurance claim incident:

\begingroup
\small
\setlength{\tabcolsep}{3pt}
\renewcommand{\arraystretch}{1.15}
\begin{longtable}{@{}>{\raggedright\arraybackslash}p{0.19\textwidth} >{\raggedright\arraybackslash}p{0.13\textwidth} >{\raggedright\arraybackslash}p{0.13\textwidth} >{\raggedright\arraybackslash}p{0.43\textwidth}@{}}
\toprule
Field & Type & Default & Description \\
\midrule
\endfirsthead
\toprule
Field & Type & Default & Description \\
\midrule
\endhead
\midrule
\endfoot
\bottomrule
\endlastfoot
\texttt{incident\_\allowbreak number} & string & required & Incident number (e.g., \#12345) \\
\texttt{reference\_\allowbreak number} & string & required & Reference ID (e.g., L240123) \\
\texttt{company\_\allowbreak name} & string & required & Trucking company name \\
division & string & General & Company division \\
\texttt{coverage\_\allowbreak type} & string & required & Coverage type (Liability, Physical Damage, Inland Marine, Cargo) \\
status & string & required & Open or Closed \\
\texttt{policy\_\allowbreak number} & string & required & Policy identifier \\
\texttt{policy\_\allowbreak state} & string & required & Policy state abbreviation \\
\texttt{cause\_\allowbreak code} & optional string & null & Internal cause code \\
description & string & required & Detailed incident description \\
handler & string & Claims Adjuster & Claims handler \\
\texttt{unit\_\allowbreak number} & optional string & null & Vehicle/truck unit ID \\
\texttt{date\_\allowbreak of\_\allowbreak loss} & string & required & Date incident occurred \\
\texttt{loss\_\allowbreak state} & string & required & State where loss occurred \\
\texttt{date\_\allowbreak reported} & string & required & Date reported to insurance \\
agency & optional string & null & Insurance agency name \\
insured & string & required & Insured party name \\
claimants & list of strings & [] & List of claimants \\
\texttt{driver\_\allowbreak name} & optional string & null & Driver name at time of incident \\
bi & financial breakdown & \{\} & Bodily Injury \\
pd & financial breakdown & \{\} & Property Damage \\
lae & financial breakdown & \{\} & Loss Adjustment Expense \\
ded & financial breakdown & \{\} & Deductible \\
\texttt{adjuster\_\allowbreak notes} & optional string & null & Additional adjuster notes \\
\end{longtable}
\endgroup

\subsection*{A.3 Extraction Output Shape}

Claim/loss-run tasks return incident records:

\begingroup
\small
\setlength{\tabcolsep}{3pt}
\renewcommand{\arraystretch}{1.15}
\begin{tabular}{@{}>{\raggedright\arraybackslash}p{0.16\textwidth} >{\raggedright\arraybackslash}p{0.20\textwidth} >{\raggedright\arraybackslash}p{0.16\textwidth} >{\raggedright\arraybackslash}p{0.36\textwidth}@{}}
\toprule
Field & Type & Default & Description \\
\midrule
incidents & list of LossRunIncident & required & List of extracted incident records \\
\bottomrule
\end{tabular}
\endgroup

Other long-list tasks return generic records:

\begingroup
\small
\setlength{\tabcolsep}{3pt}
\renewcommand{\arraystretch}{1.15}
\begin{tabular}{@{}>{\raggedright\arraybackslash}p{0.16\textwidth} >{\raggedright\arraybackslash}p{0.20\textwidth} >{\raggedright\arraybackslash}p{0.16\textwidth} >{\raggedright\arraybackslash}p{0.36\textwidth}@{}}
\toprule
Field & Type & Default & Description \\
\midrule
records & list of objects & required & List of extracted target records using the fields specified for that document family \\
\bottomrule
\end{tabular}
\endgroup

The repository publishes standalone JSON Schema files for \path{loss_run_incident}, \path{loss_run_claim_row}, \path{ifta_multisection_jurisdiction_row}, \path{policy_schedule_record}, and \path{policy_packet_item}. The multisection IFTA schema requires each jurisdiction row to include return context (\path{return_id}, year, quarter, filing type), Schedule A mileage/gallon fields, and Jurisdictions tax-detail fields.

\subsection*{A.4 Extraction Prompt Template}

Prompt-based claim multi-hop incident extraction uses the core template shown below. The evaluation code substitutes \texttt{\{schema\_json\}} with the JSON schema for the extraction object and \texttt{\{ocr\_text\}} with the selected transcript.

\begingroup
\small
\begin{verbatim}
Extract all incident records from the following document.

Requirements:
- Return ALL incidents you can find in the document.
- Each incident MUST include ALL fields in the schema.
- When a value is unknown:
  - For required string fields, use "" (empty string), not null.
  - For optional fields, use null.
  - For list fields, use [].
  - For numeric fields, use 0.0.
- Output MUST be valid JSON that conforms to the schema.

Schema (JSON Schema):
{schema_json}

Output JSON shape:
{
  "incidents": [ ... ]
}

Document:
{ocr_text}
\end{verbatim}
\endgroup

For generic operations and policy rows, the runner derives a sample-specific field contract from the keys and record groups present in the ground-truth object structure. This is schema disclosure: the contract contains field and group names but no target values, target counts, or ground-truth files. The output shape is:

\begingroup
\small
\begin{verbatim}
{
  "records": [ ... ]
}
\end{verbatim}
\endgroup

\subsection*{A.5 Agentic Sandbox Instructions}

The released Codex and Claude Code runners use the same task prompt. The temporary workspace contains only the OCR transcript, a field contract, the prompt, and an empty output directory. A macOS profile denies access to the benchmark repository, and the prompt prohibits using other host files. For claim tasks, the output path is \path{output/incidents.json}; generic tasks use \path{output/records.json}.

\begingroup
\small
\begin{verbatim}
You are dispatched as a subagent for a specific extraction task; skip
superpowers skills.

Your workspace contains:
- document_ocr.md: OCR transcript of one benchmark PDF.
- field_contract.md: the target output shape and allowed fields.

The official ground truth is not present. Do not use files outside this
workspace.

Task:
Extract the complete target list from document_ocr.md according to
field_contract.md.
Write valid JSON to output/incidents.json.

You may inspect the document, write temporary scripts, and verify counts.
Before finishing:
- Confirm the JSON parses.
- Spot-check beginning, middle, and end records against the OCR.
- Ensure field names match field_contract.md exactly.
- Ensure every target row or target record type requested by the contract is
  represented.

Do not include explanations in the JSON file.
\end{verbatim}
\endgroup

The field contract contains the claim JSON Schema or the generic record groups and allowed fields described in Appendix~A.4. It contains no record values or target counts.

\subsection*{A.6 Field Scoring Rules}

We compute strict normalized-record completeness first, then schema-conformant field precision/recall/F1 as partial-credit diagnostics.

\paragraph{Step 1: Canonicalize each incident (schema validation + normalization).}
Both predictions and ground truth are parsed as full \texttt{LossRunIncident} objects and then normalized:
\begin{itemize}
    \item \textbf{Incident identifier}: common prefixes in \path{incident_number} (e.g., \texttt{Incident \#}, \texttt{\#}, \texttt{INC}) are stripped; whitespace and case are normalized.
    \item \textbf{String fields}: all strings are trimmed and compared case-insensitively. For optional string fields, the empty string is treated as \texttt{null}. Optional string fields are:
    \begin{itemize}
        \item \texttt{cause\_code}
        \item \texttt{unit\_number}
        \item \texttt{agency}
        \item \texttt{driver\_name}
        \item \texttt{adjuster\_notes}
    \end{itemize}
    \item \textbf{Dates}: supported date forms are normalized to \texttt{MM/DD/YYYY}.
    \item \textbf{Claimants}: coerced to a list; each entry is trimmed and case-normalized; empty entries are dropped; the list is sorted.
    \item \textbf{Financial breakdowns}: for each breakdown object (\texttt{bi}, \texttt{pd}, \texttt{lae}, \texttt{ded}), we ensure it is an object and normalize each numeric field by converting to float, rounding to two decimals, and mapping negative zero to zero. Unparseable values are treated as 0.0. Zero-valued default breakdown fields are not scored. Breakdown numeric fields are:
    \begin{itemize}
        \item \texttt{reserve}
        \item \texttt{paid}
        \item \texttt{recovered}
        \item \texttt{total\_incurred}
    \end{itemize}
\end{itemize}

Generic records are canonicalized recursively. Blank fields are omitted; strings are whitespace-normalized and compared case-insensitively; dates, decimals, percentages, currency values, and accounting-parenthesis negatives are canonicalized; and list values are order-normalized. State and jurisdiction names normalize to two-letter codes, fuel descriptions to parenthetical codes, policy line-of-business names to BOP/WC/CGL acronyms, and \texttt{Applies within} clause-scope wrappers to the core scope phrase. Records are matched greedily by overlapping non-global field pairs for field diagnostics. The \path{record_type} field, hidden row identifiers, and repeated document-level policy fields are excluded from field matching and field-pair scoring. Strict exact-record comparison uses the complete normalized public record.

\paragraph{Step 2: Compare complete normalized records.}
Let $R_G$ and $R_P$ be multisets of normalized ground-truth and predicted records. Exact-record recall is $|R_G \cap R_P|/|R_G|$. A document is complete only when $R_G=R_P$, so missing records, extra records, wrong fields, and duplicate multiplicity all affect strict completeness. Record-list order is not scored in this release.

\paragraph{Step 3: Flatten incidents into a multiset of field-value pairs.}
For each canonicalized incident, we emit a list of triples (\path{incident_id}, \path{field_path}, \path{value}). Nested financial fields are represented with dotted paths (e.g., \path{bi.reserve}).

\paragraph{Step 4: Compute partial-credit field diagnostics.}
Let $\mathcal{F}(G)$ be the multiset of flattened triples from the ground truth and $\mathcal{F}(P)$ the multiset from predictions. We compute
\begin{align}
\mathrm{found} &= |\mathcal{F}(G) \cap \mathcal{F}(P)|,\\
\mathrm{recall} &= \frac{\mathrm{found}}{|\mathcal{F}(G)|},\\
\mathrm{precision} &= \frac{\mathrm{found}}{|\mathcal{F}(P)|},\\
\mathrm{F1} &= \frac{2\,\mathrm{precision}\,\mathrm{recall}}{\mathrm{precision}+\mathrm{recall}}.
\end{align}
The multiset formulation means that if a document contains exact duplicate incidents (same normalized \path{incident_id}), they are counted with multiplicity, and the score only credits matches up to the minimum count in each multiset.

\paragraph{Additional diagnostics.}
We also report:
\begin{itemize}
    \item \textbf{Missing/extra incident IDs}: computed from the set of normalized incident IDs present in each list.
    \item \textbf{Exact-record precision and $F_1$}: computed from the same normalized-record intersection when systems over-extract.
    \item \textbf{Evaluation-role and stressor slices}: aggregated from the public family and problem metadata.
\end{itemize}
